% ============================================================
% 所有表格  —  LaTeX代码
% 由 generate_tables.py 自动生成
% 使用方式: % ============================================================
% 所有表格  —  LaTeX代码
% 由 generate_tables.py 自动生成
% 使用方式: % ============================================================
% 所有表格  —  LaTeX代码
% 由 generate_tables.py 自动生成
% 使用方式: % ============================================================
% 所有表格  —  LaTeX代码
% 由 generate_tables.py 自动生成
% 使用方式: \input{all_tables.tex}  或复制对应表格到论文正文
% ============================================================

\usepackage{booktabs}
\usepackage{multirow}
\usepackage{amsmath}



\begin{table}[htbp]
  \centering
  \caption{Descriptive statistics of daily logarithmic returns.}
  \label{tab:desc_stats}
  \begin{tabular}{lc}
    \toprule
    Statistic & Value \\
    \midrule
    N & 8643 \\
    Mean (%) & 0.0396 \\
    Std. Dev. (%) & 2.0968 \\
    Median (%) & 0.0524 \\
    Min (%) & -17.9051 \\
    Max (%) & 71.9152 \\
    Q1 (%) & -0.6832 \\
    Q3 (%) & 0.7888 \\
    Skewness & 5.5063 \\
    Excess Kurtosis & 176.9019 \\
    Jarque-Bera & 11300392.8815 \\
    J-B p-value & 0.0000e+00 \\
    \bottomrule
  \end{tabular}
\end{table}

\begin{table}[htbp]
  \centering
  \caption{Unit root test results (ADF \& KPSS).}
  \label{tab:unit_root}
  \begin{tabular}{lccc}
    \toprule
    Test & Statistic & P-value & Conclusion \\
    \midrule
    ADF & -12.6088 & 0.0000 & 平稳 \\
    KPSS & 0.0552 & 0.1000 & 平稳 \\
    \bottomrule
  \end{tabular}
\end{table}

\begin{table}[htbp]
  \centering
  \caption{ARCH-LM test for conditional heteroscedasticity.}
  \label{tab:arch_test}
  \begin{tabular}{lcc}
    \toprule
    Test & Statistic & P-value \\
    \midrule
    ARCH-LM lag=10 & 610.4411 & 0.0000 \\
    \bottomrule
  \end{tabular}
\end{table}

\begin{table}[htbp]
  \centering
  \caption{Ljung-Box test for ARIMA residual white noise.}
  \label{tab:arima_resid}
  \begin{tabular}{lcc}
    \toprule
    Lag & LB Statistic & P-value \\
    \midrule
    6 & 4.0519 & 0.6696 \\
    12 & 12.8861 & 0.3774 \\
    18 & 36.0467 & 0.0070 \\
    24 & 47.5468 & 0.0029 \\
    \bottomrule
  \end{tabular}
\end{table}

\begin{table}[htbp]
  \centering
  \caption{Parameter estimation results of GARCH-family models.}
  \label{tab:model_params}
  \begin{tabular}{lcccccccccc}
    \toprule
    模型 & $\omega$ & $\alpha_1$ & $\beta_1$ & $\gamma_1$ & $\delta$ & $\nu$ & $\alpha_1+\beta_1$ & AIC & BIC & LogLik \\
    \midrule
    GARCH(1,1)-N & 0.0229 & 0.0776 & 0.9131 &  &  &  & 0.9908 & 20990.6983 & 21010.9900 & -10492.3491 \\
    GARCH(1,1)-t & 0.0177 & 0.0706 & 0.9240 &  &  & 4.7047 & 0.9946 & 20398.9667 & 20426.0223 & -10195.4834 \\
    EGARCH(1,1)-t & 0.0155 & 0.1547 & 0.9899 &  &  & 4.7706 & 0.9899 & 20379.5766 & 20406.6322 & -10185.7883 \\
    APARCH(1,1)-t & 0.0196 & 0.0750 & 0.9250 &  & 1.5532 & 4.9539 & 0.9250 & 20387.3248 & 20421.1442 & -10188.6624 \\
    \bottomrule
  \end{tabular}
\end{table}

\begin{table}[htbp]
  \centering
  \caption{Model comparison based on information criteria.}
  \label{tab:model_comp}
  \begin{tabular}{lcccc}
    \toprule
    模型 & AIC & BIC & LogLik & $\alpha_1+\beta_1$ \\
    \midrule
    GARCH(1,1)-N & 20990.6983 & 21010.9900 & -10492.3491 & 0.9908 \\
    GARCH(1,1)-t & 20398.9667 & 20426.0223 & -10195.4834 & 0.9946 \\
    EGARCH(1,1)-t & 20379.5766 & 20406.6322 & -10185.7883 & 0.9899 \\
    APARCH(1,1)-t & 20387.3248 & 20421.1442 & -10188.6624 & 0.9250 \\
    \bottomrule
  \end{tabular}
\end{table}

\begin{table}[htbp]
  \centering
  \caption{Standardized residual statistics across models.}
  \label{tab:std_resid}
  \begin{tabular}{lcccc}
    \toprule
    模型 & 均值 & 标准差 & 偏度 & 峰度(超额) \\
    \midrule
    GARCH(1,1)-N & -0.0067 & 1.0003 & -0.4095 & 4.1046 \\
    GARCH(1,1)-t & -0.0062 & 0.9924 & -0.4053 & 4.1849 \\
    EGARCH(1,1)-t & -0.0061 & 0.9954 & -0.4692 & 4.8975 \\
    APARCH(1,1)-t & -0.0062 & 1.0100 & -0.4152 & 4.2137 \\
    \bottomrule
  \end{tabular}
\end{table}

\begin{table}[htbp]
  \centering
  \caption{Rolling forecast evaluation metrics.}
  \label{tab:forecast_eval}
  \begin{tabular}{lccccc}
    \toprule
    模型 & RMSE & MAE & SMAPE(%) & QLIKE & n \\
    \midrule
    GARCH(1,1)-N & 0.8340 & 0.6151 & 82.9513 & 0.9381 & 640 \\
    GARCH(1,1)-t & 0.8404 & 0.6207 & 83.2926 & 0.9369 & 640 \\
    EGARCH(1,1)-t & 0.8181 & 0.5975 & 82.1727 & 0.9626 & 640 \\
    APARCH(1,1)-t & 0.8222 & 0.6048 & 82.6125 & 0.9289 & 640 \\
    \bottomrule
  \end{tabular}
\end{table}

\begin{table}[htbp]
  \centering
  \caption{Diebold-Mariano test for predictive accuracy comparison.}
  \label{tab:dm_test}
  \begin{tabular}{lcccc}
    \toprule
    模型1 & 模型2 & DM统计量 & DM_p值 & 结论 \\
    \midrule
    GARCH(1,1)-N & GARCH(1,1)-t & -2.4198 & 0.0158 & GARCH(1,1)-N 更优** \\
    GARCH(1,1)-N & EGARCH(1,1)-t & 1.9327 & 0.0537 & EGARCH(1,1)-t 更优* \\
    GARCH(1,1)-N & APARCH(1,1)-t & 2.3000 & 0.0218 & APARCH(1,1)-t 更优** \\
    GARCH(1,1)-t & EGARCH(1,1)-t & 3.4601 & 0.0006 & EGARCH(1,1)-t 更优*** \\
    GARCH(1,1)-t & APARCH(1,1)-t & 5.3838 & 0.0000 & APARCH(1,1)-t 更优*** \\
    EGARCH(1,1)-t & APARCH(1,1)-t & -1.2291 & 0.2195 & EGARCH(1,1)-t 略优(不显著) \\
    \bottomrule
  \end{tabular}
\end{table}

\begin{table}[htbp]
  \centering
  \caption{Summary of modeling process.}
  \label{tab:summary}
  \begin{tabular}{ll}
    \toprule
    Item & Value \\
    \midrule
    数据起始 & 2000-01-04 \\
    数据结束 & 2026-06-08 \\
    有效观测 & 6399 \\
    均值方程 & ARIMA (auto_arima选择) \\
    波动率模型数 & 4 (GARCH-N, GARCH-t, EGARCH-t, APARCH-t) \\
    样本内最优 & EGARCH(1,1)-t \\
    最优AIC & 20379.5766 \\
    样本外最优RMSE & EGARCH(1,1)-t (0.818117) \\
    \bottomrule
  \end{tabular}
\end{table}  或复制对应表格到论文正文
% ============================================================

\usepackage{booktabs}
\usepackage{multirow}
\usepackage{amsmath}



\begin{table}[htbp]
  \centering
  \caption{Descriptive statistics of daily logarithmic returns.}
  \label{tab:desc_stats}
  \begin{tabular}{lc}
    \toprule
    Statistic & Value \\
    \midrule
    N & 8643 \\
    Mean (%) & 0.0396 \\
    Std. Dev. (%) & 2.0968 \\
    Median (%) & 0.0524 \\
    Min (%) & -17.9051 \\
    Max (%) & 71.9152 \\
    Q1 (%) & -0.6832 \\
    Q3 (%) & 0.7888 \\
    Skewness & 5.5063 \\
    Excess Kurtosis & 176.9019 \\
    Jarque-Bera & 11300392.8815 \\
    J-B p-value & 0.0000e+00 \\
    \bottomrule
  \end{tabular}
\end{table}

\begin{table}[htbp]
  \centering
  \caption{Unit root test results (ADF \& KPSS).}
  \label{tab:unit_root}
  \begin{tabular}{lccc}
    \toprule
    Test & Statistic & P-value & Conclusion \\
    \midrule
    ADF & -12.6088 & 0.0000 & 平稳 \\
    KPSS & 0.0552 & 0.1000 & 平稳 \\
    \bottomrule
  \end{tabular}
\end{table}

\begin{table}[htbp]
  \centering
  \caption{ARCH-LM test for conditional heteroscedasticity.}
  \label{tab:arch_test}
  \begin{tabular}{lcc}
    \toprule
    Test & Statistic & P-value \\
    \midrule
    ARCH-LM lag=10 & 610.4411 & 0.0000 \\
    \bottomrule
  \end{tabular}
\end{table}

\begin{table}[htbp]
  \centering
  \caption{Ljung-Box test for ARIMA residual white noise.}
  \label{tab:arima_resid}
  \begin{tabular}{lcc}
    \toprule
    Lag & LB Statistic & P-value \\
    \midrule
    6 & 4.0519 & 0.6696 \\
    12 & 12.8861 & 0.3774 \\
    18 & 36.0467 & 0.0070 \\
    24 & 47.5468 & 0.0029 \\
    \bottomrule
  \end{tabular}
\end{table}

\begin{table}[htbp]
  \centering
  \caption{Parameter estimation results of GARCH-family models.}
  \label{tab:model_params}
  \begin{tabular}{lcccccccccc}
    \toprule
    模型 & $\omega$ & $\alpha_1$ & $\beta_1$ & $\gamma_1$ & $\delta$ & $\nu$ & $\alpha_1+\beta_1$ & AIC & BIC & LogLik \\
    \midrule
    GARCH(1,1)-N & 0.0229 & 0.0776 & 0.9131 &  &  &  & 0.9908 & 20990.6983 & 21010.9900 & -10492.3491 \\
    GARCH(1,1)-t & 0.0177 & 0.0706 & 0.9240 &  &  & 4.7047 & 0.9946 & 20398.9667 & 20426.0223 & -10195.4834 \\
    EGARCH(1,1)-t & 0.0155 & 0.1547 & 0.9899 &  &  & 4.7706 & 0.9899 & 20379.5766 & 20406.6322 & -10185.7883 \\
    APARCH(1,1)-t & 0.0196 & 0.0750 & 0.9250 &  & 1.5532 & 4.9539 & 0.9250 & 20387.3248 & 20421.1442 & -10188.6624 \\
    \bottomrule
  \end{tabular}
\end{table}

\begin{table}[htbp]
  \centering
  \caption{Model comparison based on information criteria.}
  \label{tab:model_comp}
  \begin{tabular}{lcccc}
    \toprule
    模型 & AIC & BIC & LogLik & $\alpha_1+\beta_1$ \\
    \midrule
    GARCH(1,1)-N & 20990.6983 & 21010.9900 & -10492.3491 & 0.9908 \\
    GARCH(1,1)-t & 20398.9667 & 20426.0223 & -10195.4834 & 0.9946 \\
    EGARCH(1,1)-t & 20379.5766 & 20406.6322 & -10185.7883 & 0.9899 \\
    APARCH(1,1)-t & 20387.3248 & 20421.1442 & -10188.6624 & 0.9250 \\
    \bottomrule
  \end{tabular}
\end{table}

\begin{table}[htbp]
  \centering
  \caption{Standardized residual statistics across models.}
  \label{tab:std_resid}
  \begin{tabular}{lcccc}
    \toprule
    模型 & 均值 & 标准差 & 偏度 & 峰度(超额) \\
    \midrule
    GARCH(1,1)-N & -0.0067 & 1.0003 & -0.4095 & 4.1046 \\
    GARCH(1,1)-t & -0.0062 & 0.9924 & -0.4053 & 4.1849 \\
    EGARCH(1,1)-t & -0.0061 & 0.9954 & -0.4692 & 4.8975 \\
    APARCH(1,1)-t & -0.0062 & 1.0100 & -0.4152 & 4.2137 \\
    \bottomrule
  \end{tabular}
\end{table}

\begin{table}[htbp]
  \centering
  \caption{Rolling forecast evaluation metrics.}
  \label{tab:forecast_eval}
  \begin{tabular}{lccccc}
    \toprule
    模型 & RMSE & MAE & SMAPE(%) & QLIKE & n \\
    \midrule
    GARCH(1,1)-N & 0.8340 & 0.6151 & 82.9513 & 0.9381 & 640 \\
    GARCH(1,1)-t & 0.8404 & 0.6207 & 83.2926 & 0.9369 & 640 \\
    EGARCH(1,1)-t & 0.8181 & 0.5975 & 82.1727 & 0.9626 & 640 \\
    APARCH(1,1)-t & 0.8222 & 0.6048 & 82.6125 & 0.9289 & 640 \\
    \bottomrule
  \end{tabular}
\end{table}

\begin{table}[htbp]
  \centering
  \caption{Diebold-Mariano test for predictive accuracy comparison.}
  \label{tab:dm_test}
  \begin{tabular}{lcccc}
    \toprule
    模型1 & 模型2 & DM统计量 & DM_p值 & 结论 \\
    \midrule
    GARCH(1,1)-N & GARCH(1,1)-t & -2.4198 & 0.0158 & GARCH(1,1)-N 更优** \\
    GARCH(1,1)-N & EGARCH(1,1)-t & 1.9327 & 0.0537 & EGARCH(1,1)-t 更优* \\
    GARCH(1,1)-N & APARCH(1,1)-t & 2.3000 & 0.0218 & APARCH(1,1)-t 更优** \\
    GARCH(1,1)-t & EGARCH(1,1)-t & 3.4601 & 0.0006 & EGARCH(1,1)-t 更优*** \\
    GARCH(1,1)-t & APARCH(1,1)-t & 5.3838 & 0.0000 & APARCH(1,1)-t 更优*** \\
    EGARCH(1,1)-t & APARCH(1,1)-t & -1.2291 & 0.2195 & EGARCH(1,1)-t 略优(不显著) \\
    \bottomrule
  \end{tabular}
\end{table}

\begin{table}[htbp]
  \centering
  \caption{Summary of modeling process.}
  \label{tab:summary}
  \begin{tabular}{ll}
    \toprule
    Item & Value \\
    \midrule
    数据起始 & 2000-01-04 \\
    数据结束 & 2026-06-08 \\
    有效观测 & 6399 \\
    均值方程 & ARIMA (auto_arima选择) \\
    波动率模型数 & 4 (GARCH-N, GARCH-t, EGARCH-t, APARCH-t) \\
    样本内最优 & EGARCH(1,1)-t \\
    最优AIC & 20379.5766 \\
    样本外最优RMSE & EGARCH(1,1)-t (0.818117) \\
    \bottomrule
  \end{tabular}
\end{table}  或复制对应表格到论文正文
% ============================================================

\usepackage{booktabs}
\usepackage{multirow}
\usepackage{amsmath}



\begin{table}[htbp]
  \centering
  \caption{Descriptive statistics of daily logarithmic returns.}
  \label{tab:desc_stats}
  \begin{tabular}{lc}
    \toprule
    Statistic & Value \\
    \midrule
    N & 8643 \\
    Mean (%) & 0.0396 \\
    Std. Dev. (%) & 2.0968 \\
    Median (%) & 0.0524 \\
    Min (%) & -17.9051 \\
    Max (%) & 71.9152 \\
    Q1 (%) & -0.6832 \\
    Q3 (%) & 0.7888 \\
    Skewness & 5.5063 \\
    Excess Kurtosis & 176.9019 \\
    Jarque-Bera & 11300392.8815 \\
    J-B p-value & 0.0000e+00 \\
    \bottomrule
  \end{tabular}
\end{table}

\begin{table}[htbp]
  \centering
  \caption{Unit root test results (ADF \& KPSS).}
  \label{tab:unit_root}
  \begin{tabular}{lccc}
    \toprule
    Test & Statistic & P-value & Conclusion \\
    \midrule
    ADF & -12.6088 & 0.0000 & 平稳 \\
    KPSS & 0.0552 & 0.1000 & 平稳 \\
    \bottomrule
  \end{tabular}
\end{table}

\begin{table}[htbp]
  \centering
  \caption{ARCH-LM test for conditional heteroscedasticity.}
  \label{tab:arch_test}
  \begin{tabular}{lcc}
    \toprule
    Test & Statistic & P-value \\
    \midrule
    ARCH-LM lag=10 & 610.4411 & 0.0000 \\
    \bottomrule
  \end{tabular}
\end{table}

\begin{table}[htbp]
  \centering
  \caption{Ljung-Box test for ARIMA residual white noise.}
  \label{tab:arima_resid}
  \begin{tabular}{lcc}
    \toprule
    Lag & LB Statistic & P-value \\
    \midrule
    6 & 4.0519 & 0.6696 \\
    12 & 12.8861 & 0.3774 \\
    18 & 36.0467 & 0.0070 \\
    24 & 47.5468 & 0.0029 \\
    \bottomrule
  \end{tabular}
\end{table}

\begin{table}[htbp]
  \centering
  \caption{Parameter estimation results of GARCH-family models.}
  \label{tab:model_params}
  \begin{tabular}{lcccccccccc}
    \toprule
    模型 & $\omega$ & $\alpha_1$ & $\beta_1$ & $\gamma_1$ & $\delta$ & $\nu$ & $\alpha_1+\beta_1$ & AIC & BIC & LogLik \\
    \midrule
    GARCH(1,1)-N & 0.0229 & 0.0776 & 0.9131 &  &  &  & 0.9908 & 20990.6983 & 21010.9900 & -10492.3491 \\
    GARCH(1,1)-t & 0.0177 & 0.0706 & 0.9240 &  &  & 4.7047 & 0.9946 & 20398.9667 & 20426.0223 & -10195.4834 \\
    EGARCH(1,1)-t & 0.0155 & 0.1547 & 0.9899 &  &  & 4.7706 & 0.9899 & 20379.5766 & 20406.6322 & -10185.7883 \\
    APARCH(1,1)-t & 0.0196 & 0.0750 & 0.9250 &  & 1.5532 & 4.9539 & 0.9250 & 20387.3248 & 20421.1442 & -10188.6624 \\
    \bottomrule
  \end{tabular}
\end{table}

\begin{table}[htbp]
  \centering
  \caption{Model comparison based on information criteria.}
  \label{tab:model_comp}
  \begin{tabular}{lcccc}
    \toprule
    模型 & AIC & BIC & LogLik & $\alpha_1+\beta_1$ \\
    \midrule
    GARCH(1,1)-N & 20990.6983 & 21010.9900 & -10492.3491 & 0.9908 \\
    GARCH(1,1)-t & 20398.9667 & 20426.0223 & -10195.4834 & 0.9946 \\
    EGARCH(1,1)-t & 20379.5766 & 20406.6322 & -10185.7883 & 0.9899 \\
    APARCH(1,1)-t & 20387.3248 & 20421.1442 & -10188.6624 & 0.9250 \\
    \bottomrule
  \end{tabular}
\end{table}

\begin{table}[htbp]
  \centering
  \caption{Standardized residual statistics across models.}
  \label{tab:std_resid}
  \begin{tabular}{lcccc}
    \toprule
    模型 & 均值 & 标准差 & 偏度 & 峰度(超额) \\
    \midrule
    GARCH(1,1)-N & -0.0067 & 1.0003 & -0.4095 & 4.1046 \\
    GARCH(1,1)-t & -0.0062 & 0.9924 & -0.4053 & 4.1849 \\
    EGARCH(1,1)-t & -0.0061 & 0.9954 & -0.4692 & 4.8975 \\
    APARCH(1,1)-t & -0.0062 & 1.0100 & -0.4152 & 4.2137 \\
    \bottomrule
  \end{tabular}
\end{table}

\begin{table}[htbp]
  \centering
  \caption{Rolling forecast evaluation metrics.}
  \label{tab:forecast_eval}
  \begin{tabular}{lccccc}
    \toprule
    模型 & RMSE & MAE & SMAPE(%) & QLIKE & n \\
    \midrule
    GARCH(1,1)-N & 0.8340 & 0.6151 & 82.9513 & 0.9381 & 640 \\
    GARCH(1,1)-t & 0.8404 & 0.6207 & 83.2926 & 0.9369 & 640 \\
    EGARCH(1,1)-t & 0.8181 & 0.5975 & 82.1727 & 0.9626 & 640 \\
    APARCH(1,1)-t & 0.8222 & 0.6048 & 82.6125 & 0.9289 & 640 \\
    \bottomrule
  \end{tabular}
\end{table}

\begin{table}[htbp]
  \centering
  \caption{Diebold-Mariano test for predictive accuracy comparison.}
  \label{tab:dm_test}
  \begin{tabular}{lcccc}
    \toprule
    模型1 & 模型2 & DM统计量 & DM_p值 & 结论 \\
    \midrule
    GARCH(1,1)-N & GARCH(1,1)-t & -2.4198 & 0.0158 & GARCH(1,1)-N 更优** \\
    GARCH(1,1)-N & EGARCH(1,1)-t & 1.9327 & 0.0537 & EGARCH(1,1)-t 更优* \\
    GARCH(1,1)-N & APARCH(1,1)-t & 2.3000 & 0.0218 & APARCH(1,1)-t 更优** \\
    GARCH(1,1)-t & EGARCH(1,1)-t & 3.4601 & 0.0006 & EGARCH(1,1)-t 更优*** \\
    GARCH(1,1)-t & APARCH(1,1)-t & 5.3838 & 0.0000 & APARCH(1,1)-t 更优*** \\
    EGARCH(1,1)-t & APARCH(1,1)-t & -1.2291 & 0.2195 & EGARCH(1,1)-t 略优(不显著) \\
    \bottomrule
  \end{tabular}
\end{table}

\begin{table}[htbp]
  \centering
  \caption{Summary of modeling process.}
  \label{tab:summary}
  \begin{tabular}{ll}
    \toprule
    Item & Value \\
    \midrule
    数据起始 & 2000-01-04 \\
    数据结束 & 2026-06-08 \\
    有效观测 & 6399 \\
    均值方程 & ARIMA (auto_arima选择) \\
    波动率模型数 & 4 (GARCH-N, GARCH-t, EGARCH-t, APARCH-t) \\
    样本内最优 & EGARCH(1,1)-t \\
    最优AIC & 20379.5766 \\
    样本外最优RMSE & EGARCH(1,1)-t (0.818117) \\
    \bottomrule
  \end{tabular}
\end{table}  或复制对应表格到论文正文
% ============================================================

\usepackage{booktabs}
\usepackage{multirow}
\usepackage{amsmath}



\begin{table}[htbp]
  \centering
  \caption{Descriptive statistics of daily logarithmic returns.}
  \label{tab:desc_stats}
  \begin{tabular}{lc}
    \toprule
    Statistic & Value \\
    \midrule
    N & 8643 \\
    Mean (%) & 0.0396 \\
    Std. Dev. (%) & 2.0968 \\
    Median (%) & 0.0524 \\
    Min (%) & -17.9051 \\
    Max (%) & 71.9152 \\
    Q1 (%) & -0.6832 \\
    Q3 (%) & 0.7888 \\
    Skewness & 5.5063 \\
    Excess Kurtosis & 176.9019 \\
    Jarque-Bera & 11300392.8815 \\
    J-B p-value & 0.0000e+00 \\
    \bottomrule
  \end{tabular}
\end{table}

\begin{table}[htbp]
  \centering
  \caption{Unit root test results (ADF \& KPSS).}
  \label{tab:unit_root}
  \begin{tabular}{lccc}
    \toprule
    Test & Statistic & P-value & Conclusion \\
    \midrule
    ADF & -12.6088 & 0.0000 & 平稳 \\
    KPSS & 0.0552 & 0.1000 & 平稳 \\
    \bottomrule
  \end{tabular}
\end{table}

\begin{table}[htbp]
  \centering
  \caption{ARCH-LM test for conditional heteroscedasticity.}
  \label{tab:arch_test}
  \begin{tabular}{lcc}
    \toprule
    Test & Statistic & P-value \\
    \midrule
    ARCH-LM lag=10 & 610.4411 & 0.0000 \\
    \bottomrule
  \end{tabular}
\end{table}

\begin{table}[htbp]
  \centering
  \caption{Ljung-Box test for ARIMA residual white noise.}
  \label{tab:arima_resid}
  \begin{tabular}{lcc}
    \toprule
    Lag & LB Statistic & P-value \\
    \midrule
    6 & 4.0519 & 0.6696 \\
    12 & 12.8861 & 0.3774 \\
    18 & 36.0467 & 0.0070 \\
    24 & 47.5468 & 0.0029 \\
    \bottomrule
  \end{tabular}
\end{table}

\begin{table}[htbp]
  \centering
  \caption{Parameter estimation results of GARCH-family models.}
  \label{tab:model_params}
  \begin{tabular}{lcccccccccc}
    \toprule
    模型 & $\omega$ & $\alpha_1$ & $\beta_1$ & $\gamma_1$ & $\delta$ & $\nu$ & $\alpha_1+\beta_1$ & AIC & BIC & LogLik \\
    \midrule
    GARCH(1,1)-N & 0.0229 & 0.0776 & 0.9131 &  &  &  & 0.9908 & 20990.6983 & 21010.9900 & -10492.3491 \\
    GARCH(1,1)-t & 0.0177 & 0.0706 & 0.9240 &  &  & 4.7047 & 0.9946 & 20398.9667 & 20426.0223 & -10195.4834 \\
    EGARCH(1,1)-t & 0.0155 & 0.1547 & 0.9899 &  &  & 4.7706 & 0.9899 & 20379.5766 & 20406.6322 & -10185.7883 \\
    APARCH(1,1)-t & 0.0196 & 0.0750 & 0.9250 &  & 1.5532 & 4.9539 & 0.9250 & 20387.3248 & 20421.1442 & -10188.6624 \\
    \bottomrule
  \end{tabular}
\end{table}

\begin{table}[htbp]
  \centering
  \caption{Model comparison based on information criteria.}
  \label{tab:model_comp}
  \begin{tabular}{lcccc}
    \toprule
    模型 & AIC & BIC & LogLik & $\alpha_1+\beta_1$ \\
    \midrule
    GARCH(1,1)-N & 20990.6983 & 21010.9900 & -10492.3491 & 0.9908 \\
    GARCH(1,1)-t & 20398.9667 & 20426.0223 & -10195.4834 & 0.9946 \\
    EGARCH(1,1)-t & 20379.5766 & 20406.6322 & -10185.7883 & 0.9899 \\
    APARCH(1,1)-t & 20387.3248 & 20421.1442 & -10188.6624 & 0.9250 \\
    \bottomrule
  \end{tabular}
\end{table}

\begin{table}[htbp]
  \centering
  \caption{Standardized residual statistics across models.}
  \label{tab:std_resid}
  \begin{tabular}{lcccc}
    \toprule
    模型 & 均值 & 标准差 & 偏度 & 峰度(超额) \\
    \midrule
    GARCH(1,1)-N & -0.0067 & 1.0003 & -0.4095 & 4.1046 \\
    GARCH(1,1)-t & -0.0062 & 0.9924 & -0.4053 & 4.1849 \\
    EGARCH(1,1)-t & -0.0061 & 0.9954 & -0.4692 & 4.8975 \\
    APARCH(1,1)-t & -0.0062 & 1.0100 & -0.4152 & 4.2137 \\
    \bottomrule
  \end{tabular}
\end{table}

\begin{table}[htbp]
  \centering
  \caption{Rolling forecast evaluation metrics.}
  \label{tab:forecast_eval}
  \begin{tabular}{lccccc}
    \toprule
    模型 & RMSE & MAE & SMAPE(%) & QLIKE & n \\
    \midrule
    GARCH(1,1)-N & 0.8340 & 0.6151 & 82.9513 & 0.9381 & 640 \\
    GARCH(1,1)-t & 0.8404 & 0.6207 & 83.2926 & 0.9369 & 640 \\
    EGARCH(1,1)-t & 0.8181 & 0.5975 & 82.1727 & 0.9626 & 640 \\
    APARCH(1,1)-t & 0.8222 & 0.6048 & 82.6125 & 0.9289 & 640 \\
    \bottomrule
  \end{tabular}
\end{table}

\begin{table}[htbp]
  \centering
  \caption{Diebold-Mariano test for predictive accuracy comparison.}
  \label{tab:dm_test}
  \begin{tabular}{lcccc}
    \toprule
    模型1 & 模型2 & DM统计量 & DM_p值 & 结论 \\
    \midrule
    GARCH(1,1)-N & GARCH(1,1)-t & -2.4198 & 0.0158 & GARCH(1,1)-N 更优** \\
    GARCH(1,1)-N & EGARCH(1,1)-t & 1.9327 & 0.0537 & EGARCH(1,1)-t 更优* \\
    GARCH(1,1)-N & APARCH(1,1)-t & 2.3000 & 0.0218 & APARCH(1,1)-t 更优** \\
    GARCH(1,1)-t & EGARCH(1,1)-t & 3.4601 & 0.0006 & EGARCH(1,1)-t 更优*** \\
    GARCH(1,1)-t & APARCH(1,1)-t & 5.3838 & 0.0000 & APARCH(1,1)-t 更优*** \\
    EGARCH(1,1)-t & APARCH(1,1)-t & -1.2291 & 0.2195 & EGARCH(1,1)-t 略优(不显著) \\
    \bottomrule
  \end{tabular}
\end{table}

\begin{table}[htbp]
  \centering
  \caption{Summary of modeling process.}
  \label{tab:summary}
  \begin{tabular}{ll}
    \toprule
    Item & Value \\
    \midrule
    数据起始 & 2000-01-04 \\
    数据结束 & 2026-06-08 \\
    有效观测 & 6399 \\
    均值方程 & ARIMA (auto_arima选择) \\
    波动率模型数 & 4 (GARCH-N, GARCH-t, EGARCH-t, APARCH-t) \\
    样本内最优 & EGARCH(1,1)-t \\
    最优AIC & 20379.5766 \\
    样本外最优RMSE & EGARCH(1,1)-t (0.818117) \\
    \bottomrule
  \end{tabular}
\end{table}